% Section 4.6: The interior product | Part: 1 | Depends on: sec:tensors, sec:metric | Status: done | Original lines: 1239-1265
%!TEX root = ../../../main.tex
\subsection{The interior product}\label{sec:interior}

For a $k$-form $\alpha$ (an alternating $k$-linear map on tangent vectors) and a vector $v$, the \emph{interior product} (contraction in the first slot) is the $(k-1)$-form
\begin{equation}\label{eq:interior}
(\iota_v\alpha)(w_1,\dots,w_{k-1})=\alpha(v,w_1,\dots,w_{k-1}).
\end{equation}
For $k=1$ this is exactly \eqref{eq:pairing}: $\iota_v\omega=\omega(v)$; for a $0$-form one puts $\iota_vf=0$.
% REORDERED (F1): the antiderivation rule eq:antider was item (f) of prop:wedgeprops; it is stated and proved here.
% REWRITTEN: the Cartan formula is a preview (the Lie derivative is defined in Section sec:lie).
The interior product satisfies
\begin{align}
\iota_v(\alpha\wedge\beta)&=\iota_v\alpha\wedge\beta+(-1)^k\alpha\wedge\iota_v\beta &&\text{(antiderivation, $\alpha\in\textstyle\bigwedge^k$),}\label{eq:antider}\\
\iota_v\iota_v&=0 &&\text{($\alpha$ alternating).}\notag
\end{align}
The second rule is immediate: $(\iota_v\iota_v\alpha)(\dots)=\alpha(v,v,\dots)=0$. Looking ahead, $\iota_X$ together with $\dd$ will
produce the Lie derivative of forms through Cartan's ``magic'' formula $\Lie_X=\dd\,\iota_X+\iota_X\,\dd$ (Section~\ref{sec:cartan}).

% PROOF ADDED BY REWRITER: in the original, eq:antider was only referred to (between Sections sec:wedgeprops and sec:interior).
\begin{proof}[Proof of \eqref{eq:antider}, step by step]
Both sides are linear in $\alpha$ and in $\beta$, and by Proposition~\ref{prop:formbasis} every form is a sum of wedges of $1$-forms.
So it suffices to take $\alpha=\alpha^1\wedge\cdots\wedge\alpha^k$ and $\beta=\alpha^{k+1}\wedge\cdots\wedge\alpha^{k+l}$ with $1$-forms $\alpha^i$.

\Step{1}{interior product of a wedge of $1$-forms} For $1$-forms $\alpha^1,\dots,\alpha^m$,
\[
\iota_v(\alpha^1\wedge\cdots\wedge\alpha^m)=\sum_{j=1}^m(-1)^{j-1}\alpha^j(v)\;\alpha^1\wedge\cdots\wedge\widehat{\alpha^j}\wedge\cdots\wedge\alpha^m
\]
(the hat means ``omit''). Indeed, by \eqref{eq:interior} and the determinant formula \eqref{eq:detformula}, the left side evaluated on
$(w_1,\dots,w_{m-1})$ is $\det[\alpha^i(u_b)]$ with $u_1=v$ and $u_{b+1}=w_b$. Expanding this determinant along its first column
$(\alpha^1(v),\dots,\alpha^m(v))$ gives $\sum_j(-1)^{j+1}\alpha^j(v)\,M_j$, where the minor $M_j$ is the determinant of $[\alpha^i(w_b)]_{i\neq j}$,
\ie, again by \eqref{eq:detformula}, the value of $\alpha^1\wedge\cdots\wedge\widehat{\alpha^j}\wedge\cdots\wedge\alpha^m$ on $(w_1,\dots,w_{m-1})$.

\Step{2}{split the sum} Apply Step~1 to $\alpha\wedge\beta=\alpha^1\wedge\cdots\wedge\alpha^{k+l}$. The terms with $j\le k$ add up to
$\bigl(\iota_v\alpha\bigr)\wedge\beta$ (Step~1 for $\alpha$ alone). In a term with $j=k+i$ the sign is
$(-1)^{k+i-1}=(-1)^k(-1)^{i-1}$, and these terms add up to $(-1)^k\alpha\wedge\iota_v\beta$ (Step~1 for $\beta$ alone).

\Step{3}{degree $0$} If $k=0$, $\alpha=f$ is a function, $\iota_vf=0$ and both sides equal $f\,\iota_v\beta$.
\end{proof}

\begin{example}[The interior product in the plane]\label{ex:interiorplane}
The $2$-form $\dd x\wedge\dd y=\dd x\otimes\dd y-\dd y\otimes\dd x$ acts on two vectors $a,b$ as a determinant:
$(\dd x\wedge\dd y)(a,b)=a^xb^y-a^yb^x$. Contract the first slot with $v$ and evaluate on a second vector $w$:
\begin{align*}
\bigl(\iota_v(\dd x\wedge\dd y)\bigr)(w)&=(\dd x\wedge\dd y)(v,w) &&\text{(definition \eqref{eq:interior})}\\
&=v^xw^y-v^yw^x &&\text{(determinant formula)}\\
&=\bigl(v^x\,\dd y-v^y\,\dd x\bigr)(w) &&\text{(since $\dd y(w)=w^y$, $\dd x(w)=w^x$).}
\end{align*}
Hence $\iota_v(\dd x\wedge\dd y)=v^x\,\dd y-v^y\,\dd x$. Two checks: (i) the antiderivation rule gives $\iota_v\dd x\wedge\dd y-\dd x\wedge\iota_v\dd y=v^x\dd y-v^y\dd x$ (same result, using $\iota_v\dd x=v^x$, $\iota_v\dd y=v^y$);
(ii) $\iota_v\iota_v(\dd x\wedge\dd y)=\iota_v(v^x\dd y-v^y\dd x)=v^xv^y-v^yv^x=0$.
\end{example}

% REORDERED (F5): the computation of iota_X(dx^dy^dz) was in Section sec:flux (original L2650-2659); it is pure
% algebra and is needed earlier (Synthesis II, Proposition prop:threelang). The flux interpretation stays in sec:flux.
\emph{Contracting the volume form of $\R^3$ (steps).} Let $\mu=\dd x\wedge\dd y\wedge\dd z$ and $X=X^1\partial_x+X^2\partial_y+X^3\partial_z$. Using the antiderivation rule on $\dd x\wedge(\dd y\wedge\dd z)$:
\begin{enumerate}
\item $\iota_X(\dd y\wedge\dd z)=X^2\dd z-X^3\dd y$ (same computation as in Example~\ref{ex:interiorplane}).
\item $\iota_X\mu=(\iota_X\dd x)\,\dd y\wedge\dd z-\dd x\wedge\iota_X(\dd y\wedge\dd z)=X^1\dd y\wedge\dd z-\dd x\wedge(X^2\dd z-X^3\dd y)$.
\item Expanding and using $-\dd x\wedge\dd z=\dd z\wedge\dd x$:
\begin{equation}\label{eq:iotavol}
\iota_X\mu=X^1\,\dd y\wedge\dd z+X^2\,\dd z\wedge\dd x+X^3\,\dd x\wedge\dd y .
\end{equation}
\end{enumerate}
\numcheck Group the same $3$-form the other way, $\mu=(\dd x\wedge\dd y)\wedge\dd z$, with $X=(1,2,3)$. By \eqref{eq:antider} with $k=2$,
$\iota_X\mu=\iota_X(\dd x\wedge\dd y)\wedge\dd z+(\dd x\wedge\dd y)\,\iota_X\dd z=(\dd y-2\,\dd x)\wedge\dd z+3\,\dd x\wedge\dd y
=\dd y\wedge\dd z+2\,\dd z\wedge\dd x+3\,\dd x\wedge\dd y$, which is \eqref{eq:iotavol} for $X=(1,2,3)$. \checkmark

The interior product is the bridge between the integration of $1$-forms (Section~\ref{sec:int1}), flux integrals (Section~\ref{sec:intk}) and the Gauss theorem (Section~\ref{sec:gauss}); its metric twin for multivectors is the left contraction of Section~\ref{sec:ops}.
